\subsection{Macaulay 模与平坦代数}

命题 9.— 设 \(\rho: A \to B\) 为 Noether 环之间的同态，M 为有限生成 A-模，N 为有限生成 B-模且在 A 上平坦。用 \({}^{a}\rho: \operatorname{Spec}(B) \longrightarrow \operatorname{Spec}(A)\) 表示与 \(\rho\) 相伴的映射。下列条件等价：

\begin{enumerate}
\def\labelenumi{(\roman{enumi})}
\item
  B-模 M ⊗A N 是 Macaulay 的；
\item
  \((\kappa(\mathfrak{p})\otimes_{\mathrm{A}}\mathrm{B})\) -模 \(\kappa(\mathfrak{p})\otimes_{\mathrm{A}}\mathrm{N}\) 对每个 \(\mathfrak{p}\in\mathrm{Supp}_{\mathrm{A}}(\mathrm{M})\) 是 Macaulay 的，且 \(\mathrm{A}_{\mathfrak{p}}\) -模 \(\mathrm{M}_{\mathfrak{p}}\) 对每个 \(\mathfrak{p}\in{}^{a}\rho(\mathrm{Supp}_{\mathrm{B}}(\mathrm{N}))\) 是 Macaulay 的；
\item
  对 \(\mathrm{Supp}_{\mathrm{B}}(\mathrm{N})\) 的每个其在 A 中的原像 \(\mathfrak{p}\) 属于 \(\mathrm{Supp}_{\mathrm{A}}(\mathrm{M})\) 的极大理想，\(\mathrm{A}_{\mathfrak{p}}\) -模 \(\mathrm{M}_{\mathfrak{p}}\) 与 \((\kappa(\mathfrak{p}) \otimes_{\mathrm{A}} \mathrm{B})\) -模 \(\kappa(\mathfrak{p}) \otimes_{\mathrm{A}} \mathrm{N}\) 都是 Macaulay 的。
\end{enumerate}

此外，若 B-模 N 是忠实平坦的，则这些条件蕴含 A-模 M 是 Macaulay 的。

设 q 为 B 的一个属于 \(M \otimes_{A} N\) 的支撑集的素理想。令 \(\mathfrak{p} = \rho^{-1}(\mathfrak{q})\)。由于模 \((\mathrm{M} \otimes_{\mathrm{A}} \mathrm{N})_{\mathfrak{q}}\) 等同于 \(M_{p} \otimes_{A_{p}} N_{q}\)，模 \(M_{p}\) 与 \(N_{q}\) 都非零，\(\kappa(\mathfrak{p}) \otimes_{A} N_{\mathfrak{q}}\) 亦然 ( \(n^{\circ} 6\) , 注)。\(A_{p}\) -模 \(N_{q}\) 同构于 \(N_{p}\) 的一个分式模，故是平坦的，并且有等式

\[
\operatorname{prof} _ {\mathrm{B} _ {\mathfrak {q}}} \left(\left(\mathrm{M} \otimes_ {\mathrm{A}} \mathrm{N}\right) _ {\mathfrak {q}}\right) = \operatorname{prof} _ {A_ {\mathfrak {p}}} \left(\mathrm{M} _ {\mathfrak {p}}\right) + \operatorname{prof} _ {\mathrm{B} _ {\mathfrak {q}}} (\kappa (\mathfrak {p}) \otimes_ {\mathrm{A}} \mathrm{N} _ {\mathfrak {q}})
\]

\[
\dim_ {B _ {q}} \left(\left(M \otimes_ {A} N\right) _ {q}\right) = \dim_ {A _ {p}} \left(M _ {p}\right) + \dim_ {B _ {q}} \left(\kappa (p) \otimes_ {A} N _ {q}\right)
\]

(§ 1, 第 6 节, 命题 11, b) 及注)，其中每一项都是 ≥ 0 的整数。考虑到 B \(_{q}\) -模 κ(p) ⊗A N \(_{q}\) 是 Macaulay 的当且仅当它作为 (κ(p) ⊗A B \(_{q}\) )-模是 Macaulay 的 (第 1 节, 例 4)，由此推出下面两个条件等价：

(α) Bq-模 (M ⊗A N)q 是 Macaulay 的；

(β) \(A_{p}\) -模 \(M_{p}\) 与 \((\kappa(\mathfrak{p})\otimes_{A}B_{\mathfrak{q}})\) -模 \(\kappa(\mathfrak{p})\otimes_{A}N_{\mathfrak{q}}\) 都是 Macaulay 的。现在来证明 (iii) 蕴含 (i)。设 n 为 B 的一个属于 \(M\otimes_{A}N\) 的支撑集的极大理想，\(\mathfrak{p}=\rho^{-1}(\mathfrak{n})\)。由前所述，有 \(\mathfrak{p}\in\operatorname{Supp}_{A}(M)\cap{}^{a}\rho(\operatorname{Supp}_{B}(N))\)；条件 (iii) 与第 6 节的注蕴含上面的条件 (β) 在 q=n 时成立。于是 \(B_{n}\) -模 \((M\otimes_{A}N)_{n}\) 是 Macaulay 的，由此得到 (i)。

蕴含关系 (ii)⇒ (iii) 是显然的；我们来证明 (i) 蕴含 (ii)。设 B-模 M ⊗\_A N 是 Macaulay 的。设 p 为 Supp\_A(M) 的一个元素。不妨设 (κ(p) ⊗\_A B)-模 κ(p) ⊗\_A N 非零，即存在 Supp\_B(N) 的一个位于 p 之上的素理想 q (第 6 节, 注)。B\_q-模 (M ⊗\_A N)\_q 是 Macaulay 的 (第 1 节, 例 3)；由前面已证的蕴含关系 (α) ⇒ (β) 及第 6 节的注推出：A\_p-模 M\_p 与 (κ(p) ⊗\_A B)-模 (κ(p) ⊗\_A N) 都是 Macaulay 的，由此得到 (ii)。

此外，若 N 在 A 上忠实平坦，则有 \(\kappa(\mathfrak{p})\otimes_{\mathrm{A}}\mathrm{N}\neq 0\)（对每个 \(\mathfrak{p}\in\operatorname{Spec}(\mathrm{A})\)），于是 \({}^{a}\rho(\operatorname{Supp}_{\mathrm{B}}(\mathrm{N}))=\operatorname{Spec}(\mathrm{A})\) (第 6 节, 注)，从而 (ii) 蕴含 M 是 Macaulay 的。

推论 1.--- 设 \(\rho: A \to B\) 为 Noether 局部环之间的一个局部同态，M 为非零的有限生成 A-模，N 为非零的有限生成 B-模且为平坦 A-模。为使 B-模 M \(\otimes_{A}N\) 是 Macaulay 的，必须且只需 A-模 M 是 Macaulay 的且 B/ \(\mathfrak{m}_{A}\)B-模 N/ \(\mathfrak{m}_{A}\)N 是 Macaulay 的。

事实上，由于 \(N/m_{A}N\) 非零，N 是忠实平坦 A-模 (I, § 3, 第 1 节, 定义 1)。

推论 2.-- 设 \(\rho: A \to B\) 为 Noether 环之间的一个同态，它使 \(B\) 成为忠实平坦 \(A\)-模，并设 \(M\) 为一个有限型 \(A\)-模。为使 \(B\)-模 \(M \otimes_{A} B\) 是 Macaulay 的，必须且只需 \(A\)-模 \(M\) 是 Macaulay 的，且 \(\kappa(\mathfrak{p}) \otimes_{A} B\) 对每个 \(\mathfrak{p} \in \operatorname{Supp}(M)\) 都是 Macaulay 环。

推论 3.--- 设 A 为 Noether 环，B 为有限平坦 A-代数，M 为有限生成的 Macaulay A-模。则 B-模 M⊗ \(_{A}\) B 是 Macaulay 的。

事实上，对每个 \(\mathfrak{p} \in \operatorname{Spec}(A)\)，环 \(\kappa(\mathfrak{p}) \otimes_{A} B\) 是有限 \(\kappa(\mathfrak{p})\)-代数，从而是 Macaulay 的 (第 5 节, 例 1)，于是应用命题 9。

推论 4.--- 设 A 为 Noether 环，J 为 A 的理想，M 为有限生成 A-模。用 \(\hat{A}\) 与 \(\hat{M}\) 分别表示 A 与 M 关于 J-进拓扑的分离完备化，用 S 表示 A 的乘性子集 \(1 + J\)。考察下列条件：

\begin{enumerate}
\def\labelenumi{(\roman{enumi})}
\item
  A-模 M 是 Macaulay 的
\item
  \(\widehat{\mathrm{A}}\) -模 \(\widehat{\mathrm{M}}\) 是 Macaulay 的；
\item
  \(\mathrm{S}^{-1}\mathrm{A}\) -模 \(\mathrm{S}^{-1}\mathrm{M}\) 是 Macaulay 的；
\item
  对每个极大理想 \(\mathfrak{m} \in \operatorname{Supp}(\mathrm{M}) \cap \mathrm{V}(\mathrm{J})\)，\(\mathrm{A}_{\mathfrak{m}}\) -模 \(\mathrm{M}_{\mathfrak{m}}\) 是 Macaulay 的；
\item
  对每个素理想 \(\mathfrak{p} \in \operatorname{Supp}(M)\)（满足 \(\mathfrak{p} + J \neq A\)），\(A_{\mathfrak{p}}\) -模 \(M_{\mathfrak{p}}\) 是 Macaulay 的且环 \(\kappa(\mathfrak{p}) \otimes_{A} \widehat{A}\) 是 Macaulay 的。
\end{enumerate}

条件 (ii) 至 (v) 等价，且都被 (i) 蕴含。当理想 J 含于 A 的根时，条件 (i) 至 (v) 等价。

我们知道 (i) 蕴含 (iii) (第 1 节, 例 3)，而当 J 含于 A 的根时，(iii) 与 (i) 相同（因为 S 的元素这时都可逆）。

环 \(\widehat{\mathrm{A}}\) 是 Noether 的 (III, § 3, 第 4 节, 命题 8)；它等同于 S \(^{-1}\) A 关于 S \(^{-1}\) J-进拓扑的完备化，而 \(\widehat{\mathrm{A}}\) -模 \(\widehat{\mathrm{M}}\) 等同于 S \(^{-1}\) M 关于 S \(^{-1}\) J-进拓扑的完备化 (III, § 3, 第 5 节, 命题 12)。因此，为证明条件 (ii) 至 (v) 的等价性，可把 \(A\) 换成 S \(^{-1}\) A，J 换成 S \(^{-1}\) J，M 换成 S \(^{-1}\) M；换言之，可设 J 含于 A 的根。这时 A-模 \(\widehat{\mathrm{A}}\) 是忠实平坦的（同上引文, 命题 9）。

显然 (v) 蕴含 (iv)，且 (iv) 蕴含 (i)。

(i)⇒(ii)：设 m 为 A 的一个极大理想；则 m 是  的一个位于 m 之上的极大理想，且  的每个极大理想都可这样得到 (III, § 3, 第 4 节, 命题 8)。环 κ(m) ⊗A  是一个域，从而是 Macaulay 环；若 A-模 M 是 Macaulay 的，则由命题 9 的蕴含关系 (iii)⇒(i)，Â-模 M 亦然。

(ii)⇒(v)：\(\hat{A}\) -模 \(\hat{M}\) 同构于 \(M\otimes_{A}\hat{A}\) (III, § 3, 第 4 节, 定理 3)；若它是 Macaulay 的，则由命题 9 的 (i)⇒(ii) 推出：\(\kappa(\mathfrak{p})\otimes_{A}\hat{A}\) 对每个 \(\mathfrak{p}\in\mathrm{Supp}(M)\) 都是 Macaulay 环，且 A-模 M 是 Macaulay 的。

命题 10.- 设 \(\rho: A \to B\) 为 Noether 环之间的一个使 B 成为平坦 A-模的同态。下列条件等价：

\begin{enumerate}
\def\labelenumi{(\roman{enumi})}
\item
  B 是 Macaulay 环；
\item
  对 B 的每个素理想 q，环 \(A_{\rho^{-1}(\mathfrak{q})}\) 与 \(\kappa(\rho^{-1}(\mathfrak{q}))\otimes_{A}B\) 都是 Macaulay 的；
\item
  对 B 的每个极大理想 n，环 \(A_{\rho^{-1}(\mathfrak{n})}\) 与 \(\kappa(\rho^{-1}(\mathfrak{n})) \otimes_{A} B\) 都是 Macaulay 的。
\end{enumerate}

此外，若 B 在 A 上忠实平坦，则这些条件蕴含 A 是 Macaulay 环。

这是命题 9 当 M = A、N = B 时的特殊情形。

推论 1.--- 在 Macaulay 环上有限的平坦代数都是 Macaulay 环。

推论 2. 设 A 为 Macaulay 环，n 为正整数；则 \(\mathrm{A}[\mathrm{X}_1,\dots ,\mathrm{X}_n]\) 与 \(\mathrm{A}[[\mathrm{X}_1,\dots ,\mathrm{X}_n]]\) 都是 Macaulay 环。

只需处理 \(n = 1\) 的情形。环 A{[}T{]} 是 Noether 的 (A, VIII, § 1, 第 4 节, 推论 1)，且对每个域 \(k\)，环 \(k[T]\) 是 Macaulay 环 (第 5 节, 例 2)；因此环 A{[}T{]} 是 Macaulay 的 (命题 10)，而环 A{[}{[}T{]}{]} 是 Macaulay 的 (命题 9 的推论 4)。

推论 3.--- Noether Macaulay 环上的每个有限生成代数都是级链的。

事实上，这样的代数是 Macaulay 环上的多项式环的商，因而是 Macaulay 环的商（推论 2），从而是级链的 (第 2 节)。

